\documentclass[10pt,a4paper]{amsart}
\usepackage[margin=19mm]{geometry}
\usepackage{amsmath,amssymb,mathtools}
\usepackage[scr=rsfs,cal=euler]{mathalfa}
\usepackage{xfrac}
\usepackage[unicode,hidelinks,bookmarks=false]{hyperref}
\newcommand{\ie}{i.e.\ }
\newcommand{\cf}{cf.\ }
\newcommand{\resp}{respectively\ }
\newcommand{\Subsection}{\subsection}
\providecommand{\nobreakdash}{\nobreak\hbox{-}}
\setlength{\emergencystretch}{2em}
\multlinegap0pt
\usepackage{mathtools}
\usepackage{amssymb}
\usepackage{float}
\usepackage[normalem]{ulem}
\newtheorem{proposition}{Proposition}
\newcommand{\bbC}{\mathbb{C}}
\newcommand{\be}{\mathbf{e}}
\newcommand{\fkg}{\mathfrak{g}}
\newcommand{\fkh}{\mathfrak{h}}
\newcommand{\gl}{\mathfrak{gl}}
\newcommand{\loc}{{\mathrm{loc}}}
\newcommand{\mD}{\mathcal{D}}
\newcommand{\reg}{\mathrm{reg}}
\newcommand{\rmH}{\mathrm{H}}
\newcommand{\sph}{{\mathrm{sph}}}
\title{Quantization of Preprojective Algebras and the Type A Rational Cherednik Algebra}
\author{Meiliang Liu, Jun Zhang, Hu Zhao}
\date{Source: arXiv:2609.15206v1 (CC BY 4.0)}
\begin{document}
\maketitle

\noindent\emph{Source and licence.} Quantization of Preprojective Algebras and the Type A Rational Cherednik Algebra. Version arXiv:2609.15206v1 (CC BY 4.0), distributed by arXiv under the Creative Commons Attribution 4.0 International licence, http://creativecommons.org/licenses/by/4.0/, as recorded in the arXiv licence metadata for this exact version and shown on the abstract page https://arxiv.org/abs/2609.15206v1. Full licence notice: \texttt{licenses/arxiv-2609.15206-CC-BY-4.0.txt}.\par\smallskip

\noindent\textit{Editorial scope.} Counted unit: the article's proposition sum (source0.tex lines 3656--3676). The article's complete original proof of it is delivered with it (source0.tex lines 3679--3935, one \texttt{proof} environment of 257 lines). No mathematics is written, repaired or re-derived: the statement and the proof are the article's own lines, in the article's own notation, and no retained line is substituted or re-flowed.  Retained uncounted as the source-local context the proof needs: lines 3435--3454; lines 3645--3653.  Nothing was omitted from any retained span, so the package carries no omission record.  No retained line contains a citation, so the wrapper carries no bibliography and no bibliography entry.  No retained line contains a figure environment, an image-inclusion command or drawing code, so no image is delivered and the package declares no asset dependency.  Editorial formatting only, in addition: an AMS \texttt{amsart} preamble that restates the article's own declarations and macros used by the retained lines, namely the environments \texttt{proposition} on separate counters, together with the standard packages those lines need.  The retained environments print in this document as Proposition 1 in order of appearance, whereas the article prints the same items with the numbering of its own full document.  The complete native source of the article as distributed in the arXiv e-print for this exact version is bundled unchanged as \texttt{sources/210399/source0.tex}.
\begin{proposition}[Localized Dunkl isomorphism]
\label{prop: dunkl embed loc neq 0}
Let \(\fkg\) be a finite-dimensional complex reductive Lie algebra with Cartan subalgebra \(\fkh\).
	Fix the root data \((R\subset\fkh^*, R^+ \subset R, W)\).
	The Dunkl embedding and its spherical restriction extend to
	\(\bbC[[\hbar]]\)-algebra isomorphisms
	\[
	\widetilde\Theta_{\hbar,c}^{\loc}:
	\rmH_{\hbar,c}(\fkh,W)^\loc
	\xrightarrow{\sim}
	\mD_\hbar(\fkh_{\reg})\rtimes W
	\]
	and
	\[
	\widetilde\Theta_{\hbar,c}^{\sph}:
	\be\rmH_{\hbar,c}(\fkh,W)^\loc\be
	\xrightarrow{\sim}
	\mD_\hbar(\fkh_{\reg})^W.
	\]
\end{proposition}

\begin{equation}\label{eq:power-sum-Cherednik-derivative}
	\Xi_{p_a}
	:=
	\frac1a\sum_{b=1}^n
	\bigl(G(p)^{-1}\bigr)_{a, b}\,
	\be\left(\sum_{i=1}^n x_i^{b-1}y_i\right)\be
	\ \in\
	\be\rmH_{\hbar,c}(n)^\loc\be.
\end{equation}

\begin{proposition}\label{prop: inverse dunkl power sum}
	Let \(\fkh \subset \gl_n (\bbC)\) be the Cartan subalgebra of diagonal matrices.
	Fix the root data \((R, R^+, W)\).
	The assignment
	\begin{equation}\label{for: inverse dunkl power sum}
		p_a\mapsto p_a\be,
		\qquad
		\delta^{-2}\mapsto\delta^{-2}\be,
		\qquad
		\partial_{p_a}\mapsto\Xi_{p_a},
		\qquad 1\leq a\leq n,
	\end{equation}
	extends uniquely to a \(\bbC[[\hbar]]\)-algebra morphism
	\[
	\Phi_{\hbar,c}:
	\mD_\hbar(\fkh_{\reg})^{S_n}
	\longrightarrow
	\be\rmH_{\hbar,c}(n)^\loc\be.
	\]
Furthermore, it is the inverse of the localized spherical Dunkl isomorphism in Proposition \ref{prop: dunkl embed loc neq 0}.
\end{proposition}

\begin{proof}
	The action of \(S_n\) on \(\fkh_{\reg}\) is free, and the quotient map
\[
\fkh_{\reg}\longrightarrow\fkh_{\reg}/S_n
\]
is finite \'etale. Hence \'etale descent for
differential operators gives
\[
\mD_\hbar(\fkh_{\reg})^{S_n}
\cong
\mD_\hbar(\fkh_{\reg}/S_n).
\]


Since
\[
\bbC[\fkh_{\reg}]^{S_n}
=
\bbC[p_1,\ldots,p_n,\delta^{-2}],
\]
	the uncompleted Weyl algebra of differential operators on
	\(\fkh_{\reg}/S_n\) is generated by
\[
p_1,\ldots,p_n,
\qquad
\delta^{-2},
\qquad
\partial_{p_1},\ldots,\partial_{p_n}.
\]
Its \(\hbar\)-adic completion
\(\mD_\hbar(\fkh_{\reg})^{S_n}\) is topologically generated by the same elements.
It therefore suffices to verify the defining relations for
this uncompleted Weyl algebra and then extend uniquely by
\(\hbar\)-adic continuity.


The generators of this uncompleted Weyl algebra
\[
p_1,\ldots,p_n,
\qquad
\delta^{-2},
\qquad
\partial_{p_1},\ldots,\partial_{p_n},
\]
 satisfy the following relations
	\begin{equation}\label{eq:power-sum-differential-relations}
	\begin{aligned}&
	\relax [p_a,p_b]=[p_a, \delta^{-2}]=0,\quad
	[\partial_{p_a},\partial_{p_b}]=0,
	\quad \delta^2 \delta^{-2} = \delta^{-2} \delta^{2} = 1,
	\\
	& [\partial_{p_a},F]
	  =\,\partial_{p_a}(F),
	  \quad
	  F\in\bbC[\fkh_{\reg}]^W.
	\end{aligned}
	\end{equation}
	To prove the multiplicative property, 
	we only need to check: for every
	\(F\in\bbC[\fkh_{\reg}]^W\),
	\begin{equation}\label{eq:Cherednik-power-sum-Weyl-relation}
	  [ \Xi_{p_a} ,F \be]
	  =
	  \partial_{p_a}(F) \be,
	  \qquad
	  [ \Xi_{p_a} ,\Xi_{p_b}]
	  =
	  0.
	\end{equation}
	
	We prove the first equality.
	Recall that  Dunkl operators are
	\[
	  D_i
	  =
	  \partial_{x_i}
	  +
	  c\sum_{j\neq i}(x_i-x_j)^{-1}(s_{ij}-1).
	\]
	Since \(F\) is \(S_n\)-invariant, it commutes with every \(s_{ij}\).
	Therefore
	\[
	\begin{aligned}
	  {}[D_i,F]
	  &=
	  [\partial_{x_i},F]
	  +
	  c\sum_{j\neq i}
	  \left[
	    (x_i-x_j)^{-1}(s_{ij}-1),
	    F
	  \right]                                                     \\
	  &=
	  \partial_{x_i}(F).
	\end{aligned}
	\]
	Using \eqref{eq:power-sum-Cherednik-derivative}, the fact that
	\(F\) commutes with \(\be\), and the preceding identity, we obtain
	\[
	\begin{aligned}
	  {}[ \Xi_{p_a} ,F\be]
	  &=
	  \frac{1}{a}\sum_{b=1}^n
	  \bigl(G(p)^{-1}\bigr)_{a, b}
	  \left[
	    \be\left(\sum_{i=1}^n x_i^{b-1}y_i\right)\be,
	    F\be
	  \right]                                                     \\
	  &=
	  \frac{1}{a}\sum_{b=1}^n
	  \bigl(G(p)^{-1}\bigr)_{a, b}
	  \be\left(
	    \sum_{i=1}^n x_i^{b-1}[y_i,F]
	  \right)\be                                                  \\
	  &=
	  \frac{1}{a}\sum_{b=1}^n
	  \bigl(G(p)^{-1}\bigr)_{a, b}
	  \be\left(
	    \sum_{i=1}^n x_i^{b-1}\partial_{x_i}(F)
	  \right)\be.
	\end{aligned}
	\]
	The chain rule gives
	\[
	\begin{aligned}
	  \sum_{i=1}^n x_i^{b-1}\partial_{x_i} (F)
	  &=
	  \sum_{i=1}^n x_i^{b-1}
	  \sum_{k=1}^n
	  kx_i^{k-1}\partial_{p_k}(F )                              \\
	  &=
	  \sum_{k=1}^n
	  k\left(\sum_{i=1}^n x_i^{b+k-2}\right)
	  \partial_{p_k}(F)                                          \\
	  &=
	  \sum_{k=1}^n
	  k p_{b+k-2}\partial_{p_k}(F)                               \\
	  &=
	  \sum_{k=1}^n
	  kG(p)_{b, k}\partial_{p_k}(F).
	\end{aligned}
	\]
	It follows that
	\[
	\begin{aligned}
	  {}[ \Xi_{p_a} ,F\be]
	  &=
	  \frac{1}{a}
	  \sum_{b,k=1}^n
	  \bigl(G(p)^{-1}\bigr)_{a, b}
	  kG(p)_{b, k}\partial_{p_k}(F)\be                             \\
	  &=
	  \frac{1}{a}
	  \sum_{k=1}^n
	  k\delta_{a, k}\partial_{p_k}(F)\be                           \\
	  &=
	  \,\partial_{p_a}(F)\be.
	\end{aligned}
	\]
	This proves the first identity in
	\eqref{eq:Cherednik-power-sum-Weyl-relation}. 
	
	For the second one,
	we reduce the problem to checking that \(\widetilde{\Theta}_{\hbar,c}^{\sph}
	\bigl( \Xi_{p_a} \bigr)
	 = \partial_{p_a}\),
	 since \(\widetilde{\Theta}_{\hbar,c}^{\sph}\) is an isomorphism, under which one derives \([\Xi_{p_a}, \Xi_{p_b}] = 0\) from \([\partial_{p_a}, \partial_{p_b}] = 0\). 
	In fact, 
	for \(0\leq r\leq n-1\), the chain rule yields
	\[
	  \sum_{i=1}^n x_i^r\partial_{x_i}
	  =
	  \sum_{a=1}^n
	  a p_{a+r-1}\partial_{p_a}.
	\]
	It gives
	\begin{equation}\label{eq:inverse-Hankel-vector-fields}
	  \partial_{p_a}
	  =
	  \frac{1}{a}\sum_{b=1}^n
	  \bigl(G(p)^{-1}\bigr)_{a, b}
	  \sum_{i=1}^n x_i^{b-1}\partial_{x_i}.
	\end{equation}
	Since \((s_{ij}-1)\be=0\), for \(0\leq r\leq n-1\) one has
	\begin{equation}\label{eq:spherical-image-power-sum-vector-field}
	  \widetilde\Theta_{\hbar,c}^{\sph}
	  \left(
	    \be\left(\sum_{i=1}^n x_i^r y_i\right)\be
	  \right)
	  =
	  \sum_{i=1}^n x_i^r\partial_{x_i}.
	\end{equation}
	Combining \eqref{eq:power-sum-Cherednik-derivative},
	\eqref{eq:inverse-Hankel-vector-fields}, and
	\eqref{eq:spherical-image-power-sum-vector-field}, we compute
	\[
	\begin{aligned}
	\widetilde{\Theta}_{\hbar,c}^{\sph}
	\bigl( \Xi_{p_a} \bigr)
	&=
	\widetilde{\Theta}_{\hbar,c}^{\sph}
	\left(
	\frac{1}{a}
	\sum_{b=1}^n
	\bigl(G(p)^{-1}\bigr)_{a, b}
	\be
	\left(
	\sum_{i=1}^n x_i^{b-1}y_i
	\right)
	\be
	\right)\\
	&=
	\frac{1}{a}
	\sum_{b=1}^n
	\bigl(G(p)^{-1}\bigr)_{a, b}
	\widetilde{\Theta}_{\hbar,c}^{\sph}
	\left(
	\be
	\left(
	\sum_{i=1}^n x_i^{b-1}y_i
	\right)
	\be
	\right)\\
	&=
	\frac{1}{a}
	\sum_{b=1}^n
	\bigl(G(p)^{-1}\bigr)_{a, b}
	\sum_{i=1}^n x_i^{b-1}\partial_{x_i}\\
	&=
	\partial_{p_a}.
	\end{aligned}
	\]
	Thus
	\eqref{for: inverse dunkl power sum} extends to an algebra morphism.
	
	
	Finally, under the standard identification
	\[
	\be\bigl(\mD_\hbar(\fkh_{\reg})\rtimes S_n\bigr)\be
	\cong
	\mD_\hbar(\fkh_{\reg})^{S_n},
	\]
	we have
	\begin{equation}\label{eq:spherical-dunkl-on-functions}
	\widetilde{\Theta}_{\hbar,c}^{\sph}(p_a\be)=p_a,
	\qquad
	\widetilde{\Theta}_{\hbar,c}^{\sph}(\delta^{-2}\be)=\delta^{-2}.
	\end{equation}
	The above calculations show
	that the composite of \(\Phi_{\hbar, c}\) defined by
	\eqref{for: inverse dunkl power sum} with
	\(\widetilde\Theta_{\hbar,c}^{\sph}\) is the identity on \(\mD_\hbar(\fkh_{\reg})^{S_n}\). Since
	\(\widetilde\Theta_{\hbar,c}^{\sph}\) is an isomorphism by
	Proposition~\ref{prop: dunkl embed loc neq 0}, the constructed morphism \(\Phi_{\hbar,c}\)
	is precisely
	\(\bigl(\widetilde\Theta_{\hbar,c}^{\sph}\bigr)^{-1}\).
\end{proof}

\end{document}
